%%% FIELD proof-generic-lower-fibers
Put \(s=\ell+3\), the number of independent source columns in an order-specific cumulant decomposition, and define
\[
 q_\beta=(0,1,\beta)^{\mathsf T},\qquad e_y=(0,0,1)^{\mathsf T}.
\]
Thus, in addition to \(p_0,p_1,\ldots,p_\ell\), the treatment and outcome disturbances have columns \(q_\beta\) and \(e_y\), respectively.  Fix coordinates on \(\operatorname{Sym}^r(\mathbb R^3)\), whose dimension is
\[
 N_r=\binom{r+2}{2},
\]
and let
\[
 V_r=\left[\operatorname{vec}_r(p_0^{\otimes r})\ \cdots\ \operatorname{vec}_r(p_\ell^{\otimes r})\ \operatorname{vec}_r(q_\beta^{\otimes r})\ \operatorname{vec}_r(e_y^{\otimes r})\right]\in\mathbb R^{N_r\times s}.
\]
Writing \(\kappa_{j,r}\) for the order-\(r\) cumulant of source \(j\), the complete cumulant-map definition and mutual independence give, in these coordinates,
\[
 \operatorname{vec}_r(C_r)=V_r\kappa_r,
 \qquad
 \kappa_r=(\kappa_{0,r},\ldots,\kappa_{\ell,r},\kappa_{t,r},\kappa_{y,r})^{\mathsf T}.
 \tag{1}
\]

We first construct a simultaneous full-row-rank loading locus.  The matrix with columns \(p_0,q_\beta,e_y\) is
\[
 \begin{pmatrix}
 1&0&0\\
 \gamma&1&0\\
 \beta\gamma&\beta&1
 \end{pmatrix},
\]
and has determinant one.  Hence \(p_0^{\otimes r},q_\beta^{\otimes r},e_y^{\otimes r}\) are linearly independent: an invertible change of basis sends them to three distinct pure-power coordinate tensors.  Moreover, the affine Veronese tensors
\[
 (1,u,v)^{\otimes r},\qquad (u,v)\in\mathbb R^2,
\]
span \(\operatorname{Sym}^r(\mathbb R^3)\).  Indeed, if a linear functional on the symmetric tensor space annihilated all of them, its associated homogeneous form \(h\) would satisfy \(h(1,u,v)=0\) for every \((u,v)\in\mathbb R^2\).  Successively viewing this as a univariate polynomial in \(v\) and then in \(u\) shows that every coefficient is zero, so \(h=0\).

Consequently, starting from the three independent tensors above, affine tensors can be adjoined one at a time outside the span already obtained.  If \(N_r\le \ell+3\), at most \(N_r-3\le\ell\) latent columns are needed.  Thus, for each \(r\in\{2,\ldots,R\}\), there is a set \(J_r\) of \(N_r\) source columns for which the determinant
\[
 D_r(\gamma,\beta,u,v)
 =\det\left(V_r[:,J_r]\right)
 \tag{2}
\]
is not the zero polynomial.  Since \(N_r\le N_R\le\ell+3\), this applies simultaneously to every retained order.  The polynomial
\[
 D_{\mathrm{load}}
 =\left(\prod_{r=2}^R D_r\right)
   \gamma\prod_{i=1}^{\ell}(u_i-\gamma)(v_i-\beta u_i)
 \tag{3}
\]
is nonzero because the real polynomial ring in the displayed structural coordinates is an integral domain and every factor is nonzero.  Its nonvanishing locus \(\mathcal G\) is therefore a nonempty real Zariski-open set.  On \(\mathcal G\), all selected matrices
\[
 W_r=V_r[:,J_r]
 \tag{4}
\]
are invertible, and all graph coefficients required to be nonzero are indeed nonzero.

We next construct a full-dimensional, genuinely real source-cumulant chart.  For each source choose \(R+1\) distinct real atoms \(x_0,\ldots,x_R\), with atoms on both sides of zero, and a probability vector \(w^0\) having every coordinate strictly positive and satisfying \(\sum_k w^0_kx_k=0\).  Such a choice is obtained, for example, by assigning sufficiently small positive masses to distinct interior atoms in \((-1,1)\) and solving for the two remaining positive masses at \(-1\) and \(1\).  The Vandermonde matrix
\[
 \mathsf V=(x_k^r)_{0\le r,k\le R}
\]
has determinant \(\prod_{k<k'}(x_{k'}-x_k)\ne0\).  Hence, after imposing \(m_0=1\) and \(m_1=0\), the remaining moments \((m_2,\ldots,m_R)\) are affine coordinates for the weights.  Strict positivity of the weights persists on an open neighborhood of the moment vector generated by \(w^0\).

For centered moments, the partition formula in the definition of the source cumulants has the triangular form
\[
 \kappa_r=m_r+P_r(m_2,\ldots,m_{r-1}),\qquad 2\le r\le R.
 \tag{5}
\]
This map has the polynomial inverse obtained recursively from (5).  It therefore sends the preceding moment neighborhood to an open neighborhood in \((\kappa_2,\ldots,\kappa_R)\).  Taking the product of these charts over the \(s\) sources gives an open set \(\mathcal K\) of independently assignable source-cumulant vectors.  Each point of \(\mathcal K\) is realized by centered, positive-variance, finite-support source laws.  Such a law is non-Gaussian because a nondegenerate Gaussian law cannot have finite support.  Taking the product of the source laws realizes their mutual independence and gives all moments finite.

Let \(\mathcal V\) be the set of legal parameter points whose structural coordinates lie in \(\mathcal G\) and whose source cumulants lie in \(\mathcal K\).  It is nonempty.  We now show both openness of its observable image and effect variation in every fiber over that image.

First fix \(\theta\in\mathcal V\).  If tensors \(C'_2,\ldots,C'_R\) are sufficiently close to \(\Phi_{\ell,R}(\theta)\), keep the structural coordinates and all unselected source cumulants fixed and, independently at each order, set
\[
 \kappa'_{J_r,r}
 =W_r^{-1}\left\{\operatorname{vec}_r(C'_r)
 -\sum_{j\notin J_r}\kappa_{j,r}\operatorname{vec}_r(p_j^{\otimes r})\right\}.
 \tag{6}
\]
Here \(p_t=q_\beta\) and \(p_y=e_y\).  Equation (4) makes every division in (6) valid.  At \(C'_r=C_r\), equation (1) shows that (6) returns the original selected cumulants.  Since there are finitely many sources and orders, continuity and openness of \(\mathcal K\) imply that all source-cumulant vectors remain in \(\mathcal K\) for every sufficiently small simultaneous tensor perturbation.  Equation (1) then shows that the perturbed tuple is attained.  Thus
\[
 \Omega_{\ell,R}:=\Phi_{\ell,R}(\mathcal V)
 \tag{7}
\]
is a nonempty Euclidean-open subset of the full cumulant-coordinate space, is contained in \(\mathcal O_{\ell,R}\), and consists of images of legal finite-support parameter points.

Now fix \(\zeta\in\Omega_{\ell,R}\) and choose \(\theta\in\mathcal V\) with \(\Phi_{\ell,R}(\theta)=\zeta\).  For real \(t\) near zero put
\[
 \beta(t)=\beta+t,qquad
 \gamma(t)=\gamma,qquad
 u_i(t)=u_i,qquad
 v_i(t)=v_i.
 \tag{8}
\]
Thus
\[
 b_i(t)=u_i-\gamma,qquad
 c_i(t)=v_i-(\beta+t)u_i.
 \tag{9}
\]
All factors in (3) remain nonzero for sufficiently small \(|t|\).  Let \(p_j(t)\) be the resulting columns and \(W_r(t)=V_r(t)[:,J_r]\).  The determinants in (2) remain nonzero near zero.  Keep the unselected cumulants fixed and define
\[
 \kappa_{J_r,r}(t)
 =W_r(t)^{-1}\left\{\operatorname{vec}_r(C_r)
 -\sum_{j\notin J_r}\kappa_{j,r}(0)
       \operatorname{vec}_r(p_j(t)^{\otimes r})\right\},
 \qquad 2\le r\le R.
 \tag{10}
\]
These functions are rational, hence continuous, in \(t\) on a neighborhood of zero.  By (1), they equal the original selected cumulants at zero and keep every complete tensor \(C_r\) exactly fixed.  Openness of \(\mathcal K\) supplies a common nonzero interval on which all cumulant vectors in (10) are realized by the centered finite-support laws constructed above.  Equations (8)--(10) therefore define a curve \(\theta(t)\in F_{\ell,R}(\zeta)\), while \(\beta(\theta(t))=\beta+t\) is nonconstant.

It remains to verify the asserted density.  A real polynomial that vanishes on the nonempty Euclidean-open set \(\Omega_{\ell,R}\) is identically zero: this follows by fixing all but one coordinate, applying the elementary one-variable fact on an interval, and inducting over the coordinates.  Hence \(\Omega_{\ell,R}\) is Zariski dense in the ambient cumulant space.  Since \(\Omega_{\ell,R}\subseteq\mathcal O_{\ell,R}\), the legal image is also Zariski dense in that space, and (7) is a Zariski-dense relatively Euclidean-open subset of \(\mathcal O_{\ell,R}\).  This proves the first assertion.

For completeness, covariance order requires a separate two-latent argument when deriving the universal lower bound of three.  The one-latent case is already supplied by the one-latent-depth theorem, and for \(\ell\ge3\) the first part applies at \(R=2\), because \(N_2=6\le\ell+3\).  Consider therefore \(\ell=2\).  For a candidate effect \(b\), put \(W_b=Y-bT\), and write \(\Sigma=(\sigma_{ab})_{a,b\in\{z,t,y\}}\) for the observed covariance.  Fix
\[
 \gamma=1,qquad u_1=2,qquad u_2=3,qquad
 \operatorname{Var}(\varepsilon_T)=1.
\]
At the base point take all five source variances equal to one, take \(c_1=c_2=1\), and take \(b=1\).  In coordinates \((Z,T,W_1)\), this gives
\[
 K_0=
 \begin{pmatrix}
 3&6&2\\
 6&15&5\\
 2&5&3
 \end{pmatrix},
\qquad
 \Sigma_0=
 \begin{pmatrix}
 3&6&8\\
 6&15&20\\
 8&20&28
 \end{pmatrix}.
 \tag{11}
\]
The first matrix has positive leading principal minors \(3,9,12\), so both are positive definite because the change from \((Z,T,W_1)\) to \((Z,T,Y)\) is invertible.

For \(\Sigma\) near \(\Sigma_0\), define \(A_0,A_1,A_2\) by
\[
 \begin{pmatrix}
 1&1&1\\
 1&2&3\\
 1&4&9
 \end{pmatrix}
 \begin{pmatrix}A_0\\A_1\\A_2\end{pmatrix}
 =
 \begin{pmatrix}\sigma_{zz}\\\sigma_{zt}\\\sigma_{tt}-1\end{pmatrix}.
 \tag{12}
\]
The coefficient matrix is Vandermonde and invertible, and (12) gives \(A_0=A_1=A_2=1\) at \(\Sigma_0\).  For \((\Sigma,b)\) near \((\Sigma_0,1)\), define \(d_1,d_2\) and \(A_y\) by
\[
 \begin{pmatrix}A_1&A_2\\2A_1&3A_2\end{pmatrix}
 \begin{pmatrix}d_1\\d_2\end{pmatrix}
 =
 \begin{pmatrix}
 \sigma_{zy}-b\sigma_{zt}\\
 \sigma_{ty}-b\sigma_{tt}
 \end{pmatrix},
 \tag{13}
\]
\[
 A_y=\sigma_{yy}-2b\sigma_{ty}+b^2\sigma_{tt}
       -A_1d_1^2-A_2d_2^2.
 \tag{14}
\]
At the base point, (13)--(14) give \(d_1=d_2=A_y=1\).  Therefore there are neighborhoods on which all \(A_0,A_1,A_2,A_y\) are positive, the determinant \(A_1A_2\) of the matrix in (13) is nonzero, and both \(d_i\) are nonzero.

Equations (12)--(14) give the exact covariance decomposition
\[
 \operatorname{Cov}(Z,T,W_b)=
 A_0(1,1,0)^{\otimes2}
 +A_1(1,2,d_1)^{\otimes2}
 +A_2(1,3,d_2)^{\otimes2}
 +(0,1,0)^{\otimes2}
 +A_y(0,0,1)^{\otimes2}.
 \tag{15}
\]
Indeed, its first three entries are (12), its two cross-entries with \(W_b\) are (13), and its last diagonal entry is (14).  Returning to \(Y=W_b+bT\), equation (15) is precisely the graph cumulant formula with effect \(\beta=b\), latent coefficients \(b_i=u_i-\gamma\in\{1,2\}\), and \(c_i=d_i\ne0\).  Independent centered two-point laws with the five displayed positive variances realize (15); they are finite-support and non-Gaussian.  Shrinking the neighborhoods if necessary, every covariance \(\Sigma\) in an open ball around \(\Sigma_0\) admits (15) for every \(b\) in a nontrivial interval around one.  Hence this open ball lies in \(\mathcal O_{2,2}\), and every one of its fibers contains unequal-effect legal points.  The same polynomial argument used after (10) makes this ball Zariski dense.  Thus order two cannot generically identify the effect when \(\ell=2\).

Let
\[
 R_0(\ell)=\min\left\{R\ge2:\binom{R+2}{2}>\ell+3\right\}.
\]
For every \(2\le R<R_0(\ell)\), the first part produces a Zariski-dense subset of legal observations with nonconstant effect fibers.  If such an \(R\) belonged to the identifying set in the definition of \(r_\beta(\ell)\), its nonempty relative Zariski-open singleton-fiber locus would intersect this dense nonidentification subset, a contradiction.  The preceding covariance arguments give the additional lower bound \(r_\beta(\ell)\ge3\) for every \(\ell\ge1\).  The strict upper-and-recovery theorem gives finiteness and \(r_\beta(\ell)\le U_\ell\).  Therefore
\[
 L_\ell=\max\{3,R_0(\ell)\}
 \le r_\beta(\ell)\le U_\ell.
 \tag{16}
\]

Finally, the positive roots of the defining quadratic inequalities give
\[
 d_\ell=\frac{-3+\sqrt{8\ell+9}}{2}+O(1)
          =\sqrt{2\ell}+O(1),
 \qquad
 R_0(\ell)=\frac{-3+\sqrt{8\ell+25}}{2}+O(1)
          =\sqrt{2\ell}+O(1).
 \tag{17}
\]
Also, since \(\binom{h+3}{3}\ge h^3/6\), the choice \(h=\lceil(6(\ell+1))^{1/3}\rceil\) proves \(h_\ell=O(\ell^{1/3})\).  Hence
\[
 L_\ell=\sqrt{2\ell}+O(1),
 \qquad
 U_\ell=d_\ell+h_\ell+1
       =\sqrt{2\ell}+O(\ell^{1/3}).
 \tag{18}
\]
Dividing (16) by \(\sqrt{2\ell}\) and using (18) yields
\[
 \frac{r_\beta(\ell)}{\sqrt{2\ell}}\longrightarrow1.
\]
